\section{Course Structure}
\label{sec:1.5}

This textbook grew out of a 9-week course on Monte Carlo methods that has been taught regularly
at the University of Chicago since 2019. The prerequisites for the course are multivariate
calculus, linear algebra, a first course in probability and statistics (including Markov chains),
and elementary knowledge of ordinary differential equations. The course follows the structure of
this textbook, covering the material chapter by chapter. To review Markov chain theory in
preparation, students are encouraged to read Appendix~A and consult the standard references on
Markov chains and stochastic processes cited there.

Course assessment is based on weekly homework assignments, an independent project, and a final
exam. Exercises representative of the weekly assignments appear at the end of each chapter.
In addition, students complete a 10-page independent project, which is intended to encourage
students to explore a topic of their choice more deeply, with the flexibility to focus on theory,
methodology, or applications. The final exam provides students with an opportunity to review the
core concepts and algorithms presented in this textbook.

Numerous variations of this course structure and grading schemes are possible. For instance,
a semester-long course could include a review of Markov chains before introducing the
Metropolis Hastings algorithm in Chapter~4, and the assessment could incorporate hands-on group
projects and oral presentations.
